\chapter{Experiments with Black-Scholes dynamics}
\label{sec:bs_experiments}

The Black--Scholes framework rests on: (i) a frictionless, arbitrage-free market with continuous trading and no transaction costs; (ii) a risky asset $S_t$ with \emph{constant} instantaneous drift $\mu$ and volatility $\sigma>0$; (iii) a constant risk-free rate $r$; (iv) log-normally distributed asset prices, driven by a single standard Brownian motion $W_t$ under the physical measure $\mathbb{P}$; and (v) no dividends, jumps, or stochastic volatility.

Under $\mathbb{P}$, the asset price follows a geometric Brownian motion (GBM):
\begin{equation}
dS_t = \mu S_t\, dt + \sigma S_t\, dW_t, \qquad S_0 > 0.
\end{equation}
Under the risk-neutral measure $\mathbb{Q}$, obtained via Girsanov's theorem with market price of risk $\lambda = (\mu-r)/\sigma$, the drift is replaced by $r$:
\begin{equation}
dS_t = r S_t\, dt + \sigma S_t\, dW_t^{\mathbb{Q}}.
\end{equation}


Applying It\^o's lemma to $\ln S_t$ yields the exact solution
\begin{equation}
S_t = S_0 \exp\left[\left(r - \tfrac{1}{2}\sigma^2\right)t + \sigma W_t^{\mathbb{Q}}\right],
\end{equation}
so that $\ln S_t \mid S_0 \sim \mathcal{N}\!\left(\ln S_0 + (r-\tfrac12\sigma^2)t,\; \sigma^2 t\right)$: log-returns are Gaussian and i.i.d.\ over disjoint intervals, and $S_t$ is log-normal.


Sample paths are continuous (no jumps), Markovian, and self-similar in the sense $S_{t+\Delta}/S_t \stackrel{d}{=} S_\Delta/S_0$ with instantaneous variance constant.


\section{Data and Training procedure}

For each trajectory $j = 1, ..., J$ we define 
the SDE of the log-price $X_t^j$ to be modelled using the Black-Scholes framework where the only source of uncertainty is given by the Brownian motion $W_t$:
\begin{equation}
    \label{log-price SDE}
    \begin{aligned}
       &dX_t^j = -\frac{1}{2}\sigma_j^2 dt + \sigma_j dW_t\\
         &\text{with } X_0^j = \log(S_0^j)
    \end{aligned}
    \end{equation} 

Where $S_0^j$ is the raw price at time $0$ of trajectory $j$, $\sigma_j$ is the volatility of trajectory $j$ and $W_t$ is a standard Brownian motion.
In this setting we are emulating the scenario of working under the risk-neutral measure $\mathbb{Q}$ with zero interest rate and zero dividend yield.\\
$X_0^j$ and $\sigma_j$ are randomly sampled for each trajectory such that 
$X_0^j \in [-1, 1]$ and $\sigma_j \in (0.0, 1.2]$.\\

In this context our condition $c_j$ for a general data observation $j$ is given by
a vector containing the trajectory $c_j=[X_0^j, X_{dt}^j, X_{2dt}^j, ..., X_{Ndt}^j]$.

 The SDE is discretized over a uniform time space $N dt = T$ obtaining the following form
 adopted for the simulation of the trajectories:
 \begin{equation}
    \label{eq: discretized_SDE}
    X_{t+dt}^j = X_t^j -\frac{1}{2}\sigma_j^2 dt + \sigma_j \sqrt{dt} Z_t
 \end{equation}
 Where $Z_t \sim \mathcal{N}(0,1)$.\\
The procedure to build the condition data matrix $C$ is summarized in Algorithm \ref{alg:y condition}.

\begin{algorithm}[H]
\caption{Algorithm for the condition data matrix $C$ creation}
\label{alg:y condition}
\begin{algorithmic}[1]

\STATE Define the number of simulation $J$.
\STATE Initialize an empty data matrix $C$ of shape $J\times N$.
\STATE Set iteration counter $j \leftarrow 0$.
\REPEAT
    \STATE Randomly sample $\sigma_j \in (0.1, 1.2)$ with $X_0^j = 0.0$.
    \STATE Evolve $X_t^j$ based on the discretized SDE for $0, dt, ..., Ndt$ steps.
    \STATE Store the trajectory $[X_0^j, X_{dt}^j, X_{2dt}^j, ..., X_{Ndt}^j]$ in the $j$-th row of $C$.
    \STATE Increment counter $j \leftarrow j + 1$.
\UNTIL{$j = J$}
\RETURN $C$
\end{algorithmic}
\end{algorithm}


Since the ForGAN architecture is adopted, which is trained to generate the next value of the trajectory,
the output data $X$ is a $J \times 1$ vector where each element is the $n$-th step ahead value of the trajectory, 
so the output is a single value and not a probability distribution during training time. 
In this case, the target output data matrix $X$ is computed by evolving each trajectory for the $\tau$-th step using the discretized SDE 
described in  \ref{eq: discretized_SDE} and storing the resulting value in the $j$-th position of $X$.\\

Now that the training data are defined the ForGAN model can be trained by using both the classical BCE loss, using algorithm \ref{alg: cgan}
and the Wasserstein loss with gradient penalty using algorithm \ref{alg:cwgan_gp}.
 




After the training phase, the ForGAN is able to generate the future distribution by repeatedly sampling the next value of the trajectory by using different noise realizations, as described in \ref{alg: forgan-inference}.\\
At the end of the condition trajectory, so at time $T=N dt$, we have the following information: the final log-price $X_{Ndt}^j$ of each trajectory $j$ and the corresponding volatility $\sigma_j$.
The goal of the Conditional GAN is to learn the distribution of final values of the stochastic process at next time horizon $T + \tau$. 
A theoretical closed form can be used as ground truth. It is given by a normal distribution:
\begin{equation}
    \begin{aligned}
    X_{Ndt + \tau}^j & \quad \sim \quad \mathcal{N}(X_{Ndt}^j - \frac{1}{2}\sigma_j^2 \Delta t, \sigma_j \sqrt{\Delta t})\\
    \text{where } &\Delta t = dt \times \tau \text{ with $\tau$ the number of days ahead. }
    \end{aligned}
\end{equation}


In order to provide a comparison between the proposed ForGAN architectures and 
the theoretical solution of the GBM processes, the target distribution $x$ of the next value $c_{T + \tau}$ of the process is computed by using the Monte Carlo algorithm described in \ref{alg:mc_terminal_dist}.
This is required to properly evaluate the model's performance because since the ForGAN model provides the distribution of the next value by repeated sampling, the vector of those samplings needs to be converted into a probability distribution. 
So, it is required a uniform way to compute the bins values used to build the empirical distribution.\\
In fact, after the Monte Carlo simulations are performed, the resulting terminal values are used to build the empirical distribution of the next value $c_{T + \tau}$ by using $\sqrt{n_{mc\_simulations}}$ as rule of thumb to compute the optimal number of bins for the histogram. 
The same bins are then used to compute the empirical distribution of the generated samples from the ForGAN model.\\
Alternatively, the Freedman-Diaconis rule can be used \cite{forgan:freedman_bins} to compute the optimal bin width for complex target distributions.
The Freedman-Diaconis rule provides a robust statistical method for determining optimal histogram bin width by minimizing the integrated mean squared error. 
Unlike methods sensitive to outliers, it utilizes the interquartile range (IQR) to account for data variability. 
The bin width $h$ is calculated as:$h = 2 \cdot \frac{\text{IQR}(x)}{\sqrt[3]{n}}$ where $n$ represents the number of observations.
By prioritizing the spread of the central 50\% of the distribution, 
the algorithm ensures that the resulting binning structure accurately reflects the underlying density of the empirical data without over-smoothing or producing noise-induced artifacts.\\

The algorithm to compute the theoretical distribution $p(x|c)$ given Monte Carlo simulations and the ForGAN's generated distribution are as follows:
\begin{algorithm}[H]
\caption{Monte Carlo generation of terminal distributions}
\label{alg:mc_terminal_dist}
\begin{algorithmic}[1]
\REQUIRE Number of trajectories $J$, steps ahead $N$, and MC simulations $K$, final values from which to compute the distribution $\mathbf{C}_T \in \mathbb{R}^J$, drift $\boldsymbol{\mu} \in \mathbb{R}^J$, volatility $\boldsymbol{\sigma} \in \mathbb{R}^J$, and time increment $ dt$.
\STATE \COMMENT{Generate random standard normal shocks for all trajectory $j$ and simulation $k$}
\STATE $\mathbf{Z} \leftarrow \text{Tensor}(J, K)$ where $Z_{j, k} \sim \mathcal{N}(0,1)$
\STATE $\Delta t = dt * N$
\FOR{each trajectory $j \in \{1, \dots, J\}$}
    \STATE \COMMENT{Calculate the deterministic drift component per step}
    \STATE $\mathbf{d_j} \leftarrow (\boldsymbol{\mu_j} - 0.5 \boldsymbol{\sigma_j}^2) \Delta t$ 
    \STATE \COMMENT{Compute stochastic shocks for simulation $k$}
    \STATE $\mathbf{S}_{k}^j \leftarrow \sigma_j \cdot Z_{j, k} \cdot \sqrt{\Delta t}$
    \STATE \COMMENT{Calculate increments for each simulation $k$}
    \STATE $\boldsymbol{\delta}_{k}^j \leftarrow d_j + S_{k}^j$ 
    \FOR{each simulation $k \in \{1, \dots, K\}$}
        \STATE $X_{k}^j \leftarrow C_{T,j} + \delta_{k}^j$
    \ENDFOR
\ENDFOR
\STATE $\mathbf{X} \leftarrow \text{Matrix}(J, K)$
\STATE Compute $\mathbf{D} \leftarrow \quad$ the normalized histogram of $\mathbf{X}$ to estimate the empirical distribution for each trajectory $j$ using $n_{bins} = \sqrt{K}$ or alternatively the method in \cite{forgan:freedman_bins}.
\RETURN $\mathbf{D}$, $\mathbf{bins}$ \COMMENT{The empirical probability distributions for each trajectory}
\end{algorithmic}
\end{algorithm}






\begin{algorithm}[H]
\caption{ForGAN distribution inference}
\label{alg: forgan-inference}
\begin{algorithmic}[1]
\REQUIRE Trained generator parameters $\theta_G$, conditioning history $c_j$, number of samples $V$
\STATE Initialize an empty list $\mathcal{S} \leftarrow \emptyset$ to store generated samples.
\FOR{$i = 1$ \TO $V$}
    \STATE Sample latent noise $z^{(i)} \sim p_z$ \COMMENT{Independent noise draw, same $c_t$}
    \STATE $\hat{x}_{t+h}^{(i)} \leftarrow G(z^{(i)}, c_j)$ \COMMENT{Generate candidate forecast}
    \STATE $\mathcal{S} \leftarrow \mathcal{S} \cup \{\hat{x}_{t+h}^{(i)}\}$
\ENDFOR
\RETURN $\mathcal{S}$ \COMMENT{ list of generated samples for empirical distribution of trajectory $j$}
\end{algorithmic}
\end{algorithm}

After having obtained $D$ from \ref{alg:mc_terminal_dist} and having computed the ForGAN inferred distribution by normalizing the histogram of $\mathcal{S}$ obtained from \ref{alg: forgan-inference} by using the same $bins$ from \ref{alg:mc_terminal_dist},
the two distributions can be compared by using all the metrics described in \ref{sec:performance_metrics} like Jensen-Shannon divergence, Wasserstein distance, Total Variation distance and statistical tests like the Kolmogorov-Smirnov test.


\section{ForGAN results with variable volatility trajectories}
\label{sec:bs experiment}
The following setup for the experiments has been used:\\
\begin{itemize}
    \item \textbf{Initial price range}: $X_0 = 0.0$
    \item \textbf{Drift range}: $\mu = 0.0$
    \item \textbf{Volatility range}: $\sigma \in (0.1, 1.2)$
    \item \textbf{Time horizon}: $T = 1.0 $ yearly horizon
    \item \textbf{Number of time steps}: $N = 252$
    \item \textbf{Number of simulated paths}: $J = 100{,}000$
    \item \textbf{n-steps ahead}: $\tau = 1$
\end{itemize}

The experimental setup is based on a data generation process designed to emulate realistic financial price dynamics exploiting Geometric Brownian motion described in the before in previous section. 
The asset price trajectories are simulated over a fixed time horizon of one year, discretized into $N=252$ time steps corresponding to daily observations.
For each experiment, $J=100{,}000$ independent paths are generated, ensuring sufficient sample diversity and statistical robustness.
The initial asset price $X_0$ is set for each trajectory to $0.0$, while the drift parameter $\mu$ is fixed at zero to isolate the model's ability to learn volatility-driven dynamics.
The volatility parameter $\sigma$ is sampled from a wide range $(0.1, 1.2)$ to capture varying market conditions, from near-deterministic to highly stochastic regimes.
The distribution to be learned refers to the $1$-th step ahead observation, that is, the distribution of asset prices at time $T + 1$ (or $N dt + 1$) given information up to time $T$. 
Learning distributions for longer horizons does not provide additional insights since the shape of the distribution remains Gaussian, with the only difference being a larger variance.\\

The experiment has been conducted by using the following architecture and hyperparameters:
\begin{table}[H]
\centering
\caption{ForGAN Architecture and Hyperparameters Summary}
\label{tab:forgan_summary}
\begin{tabular}{lll}
\hline
\textbf{Parameter}         & \textbf{Generator}          & \textbf{Critic (Discriminator)} \\ \hline
Input Dimension            & 64 (Latent) + 252 (Cond.)   & 1 (Data) + 252 (Cond.)          \\
Hidden Layers              & [512, 256, 128, 64]         & [512, 256, 256, 128, 64]        \\
Output Dimension           & 1                           & 1                               \\
Activation Function        & ReLU                        & Leaky ReLU                      \\
Normalization              & Batch Norm                  & None                      \\
Dropout                    & 0.0                         & 0.0                             \\
\multicolumn{3}{c}{\textbf{Global Training Settings}}                                     \\ \hline
Learning Rate              & \multicolumn{2}{l}{0.0005}                                      \\
Max Epochs                 & \multicolumn{2}{l}{100}                                       \\
\end{tabular}
\end{table}

The Kolmogorov-Smirnov test performed on 100000 real and generated probability distributions shows that
the resulting p-values are accepted 25\% of time.\\
Moreover, the average Jensen-Shannon divergence over the test samples is around 0.49,
suggesting poor performance since the Jensen-Shannon divergence computed by using the logarithm in base 2
is bounded in the range $[0, 1]$.

\begin{table}[h]
\centering
\begin{tabular}{l c}
\hline
\textbf{Metric} & \textbf{Value} \\
\hline
KS Test Acceptance (\%) & 25.00 \\
EMD Distance & 0.00855 \\
JS Distance & 0.4378 \\
Hellinger Distance & 0.4001 \\
TV Distance & 0.4129 \\
\hline
\end{tabular}
\caption{ForGAN performance metrics averaged over 100000 samples}
\end{table}

These results strongly suggest that it is worth trying to improve the performances of the proposed ForGAN by adopting the Wasserstein distance instead of the BCE loss.\\
In order to visually assess the capabilities of the model,
different samples have been drawn such that every sampling parameter is fixed except for the volatility $\sigma$.
In particular, six different samples are shown below where
the starting values is the same $X_0 = 0.0$, $\mu = 0.0$ while the volatility have different values for each sample 
$\sigma = [0.3, 0.5, 0.7, 0.9]$. From the results can be notices that also this new architecture
struggles to effectively capture the variance of the target distribution when the volatility parameter is high.

\begin{figure}[H]
    \centering
    \includegraphics[width=1.0\textwidth]{Background/images/samples_forgan_simple.png}
    \caption{Samples obtained using fixed $X_0$ and $\mu$ while varying $\sigma = [0.3, 0.5, 0.7, 0.9]$.}
    \label{forgan_simple_sample_sigma}
\end{figure}



In case of the ForWGAN architecture experiments 
just the training procedure and the architecture differs with respect to the ForGAN case, 
in fact, both training and test data remain the same.
The experiment has been conducted by using the following architecture and hyperparameters:
\begin{table}[H]
\centering
\caption{ForWGAN Architecture and Hyperparameters Summary}
\label{tab:forwgan_summary}
\begin{tabular}{lll}
\hline
\textbf{Parameter}         & \textbf{Generator}          & \textbf{Critic (Discriminator)} \\ \hline
Input Dimension            & 64 (Latent) + 252 (Cond.)   & 1 (Data) + 252 (Cond.)          \\
Hidden Layers              & [512, 256, 128, 64]         & [512, 256, 256, 128, 64]        \\
Output Dimension           & 1                           & 1                               \\
Activation Function        & ReLU                        & Leaky ReLU                      \\
Normalization              & None                        & Layer Norm                      \\
Dropout                    & 0.0                         & 0.0                             \\
\multicolumn{3}{c}{\textbf{Global Training Settings}}                                     \\ \hline
Max Epochs                 & \multicolumn{2}{l}{100}                                       \\
Learning Rate              & \multicolumn{2}{l}{0.0005}                                      \\
\end{tabular}
\end{table}


The Kolmogorov-Smirnov test performed on 100000 real and generated probability distributions shows that
the resulting p-values are accepted 34\% of times, quite higher than the ForGAN proposed models but still far from the desired level.\\
Moreover, the average Jensen-Shannon divergence over the test samples is around 0.28 which highlights an improvement with respect to the ForGAN model.\\
\begin{table}[h]
\centering
\begin{tabular}{l c}
\hline
\textbf{Metric} & \textbf{Value} \\
\hline
KS Test Acceptance (\%) & 34.00 \\
EMD distance & 0.00380 \\
JS Distance & 0.2802 \\
Hellinger Distance & 0.2510 \\
TV Distance & 0.2539 \\
\hline
\end{tabular}
\caption{ForWGAN performance metrics averaged over 100000 samples}
\end{table}


A visual inspection of different samples with various degrees of volatility is presented below
In particular, six different samples are shown below where
the starting values is the same $X_0 = 0.0$, $\mu = 0.0$ while the volatility have different values for each sample 
$\sigma = [0.3, 0.5, 0.7, 0.9]$.
From the results we can see an important improvement with respect to the simple ForGAN architecture,
suggesting that the Wasserstein loss with gradient penalty helps in stabilizing the training process 
and improving the quality of the generated distributions the same way has been demonstrated in papers like \cite{background:wgan} and \cite{background2:gp_wgan}.

\begin{figure}[H]
    \centering
    \includegraphics[width=1.0\textwidth]{Background/images/samples_forgan.png}
    \caption{Samples obtained using fixed $X_0$ and $\mu$ while varying $\sigma = [0.3, 0.5, 0.7, 0.9]$.}
    \label{forgan_sample_sigma}
\end{figure}


